\documentclass[multi=tikzpicture,border=7pt]{standalone}
\usepackage{ctex,amsmath,bm}
\usepackage{tikz}
\usetikzlibrary{arrows.meta,calc,angles,quotes,decorations.pathmorphing,patterns,positioning,intersections,3d}
\tikzset{
axis/.style={->,>=Stealth,thick},force/.style={->,>=Stealth,very thick},
flow/.style={->,>=Stealth,thick,blue!70!black},dim/.style={<->,>=Stealth,thin},
body/.style={draw=black,fill=blue!7,thick},guide/.style={text=black,dashed,gray!70,text=black},
point/.style={circle,fill=black,inner sep=1.5pt},impulse/.style={->,>=Stealth,very thick,red!75!black}}
\begin{document}

% 01 question-01
\begin{tikzpicture}[x=1cm,y=1cm]
  \draw[thick] (-4.7,0)--(4.8,0);
  \draw[axis] (0,0)--(4.9,0) node[right]{$x$};
  \draw[axis] (0,0)--(0,4.7) node[above]{$y$};
  \coordinate (C) at (0,2);
  \draw[body] (C) circle(2);
  \foreach \a in {0,30,...,330}{\draw[thick] (C)--++(\a:2);}
  \fill (C) circle(2pt);
  \draw[dim] (C)--(0,0) node[midway,right]{$r$};
  \draw[flow,purple!70!black] (-1.0,3.1) arc[start angle=135,end angle=35,radius=1.45]
    node[midway,above]{$\omega$};
  \draw[guide] (-2.3,{2*(1-.3)})--(2.4,{2*(1-.3)});
  \filldraw[fill=orange!35,thick] (1.05,{2*(1-.3)}) circle(.11)
    node[above right]{$P\bigl(x,r(1-\lambda)\bigr)$};
  \node[red!75!black,right] at (1.2,1.12){$\odot$ 箭沿 $z$ 穿过轮面};
  \draw[flow] (0,2)--(1.4,2) node[above right]{$v_C=\omega r$};
  \node at (0,-.55){轮缘与 $12$ 根等角间隔轮辐随轮心共同平移};
\end{tikzpicture}

% 02 solution-01
\begin{tikzpicture}[x=2.2cm,y=2.7cm]
 \draw[axis] (-1.45,0)--(1.5,0) node[right]{$u$};
 \draw[axis] (0,-.7)--(0,.78) node[above]{$G_\lambda(u)$};
 \draw[red!75!black,very thick,domain=-1.266103:1.266103,samples=161] plot(\x,{\x-.3*tan(\x*180/pi)});
 \draw[blue!75!black,very thick,domain=-.927295:.927295,samples=121] plot(\x,{\x-.6*tan(\x*180/pi)});
 \draw[guide] (-1.35,{pi/12})--(1.35,{pi/12}) node[right]{$\Delta/2$};
 \draw[guide] (-1.35,{-pi/12})--(1.35,{-pi/12});
 \fill[red!75!black] (.991157,.532899) circle(.7pt) node[above right]{$g(0.3)$};
 \draw[guide] (.991157,0)--(.991157,.532899);
 \node[red!75!black,anchor=west] at (-1.38,1.05){$\lambda=0.3$：内部极值决定扫掠边界};
 \node[blue!75!black] at (.35,-.55){$\lambda=0.6$：$2g<\Delta$，有永久安全点};
 \begin{scope}[shift={(0,-1.35)},x=2.7cm,y=1cm]
  \draw[axis] (-1.08,0)--(1.1,0) node[right]{$\theta=\omega t$};
  \draw[very thick,green!50!black] (-.933939,.13)--(.973939,.13);
  \foreach \v/\lab in {-.933939/进入轮缘,-0.590366828/轮辐,-0.311222205/轮辐,0.028574763/轮辐,.973939/离开轮缘}{
   \draw[thick] (\v,-.1)--(\v,.23);
  }
  \node[below] at (-.94,-.1){进入};\node[below] at (.97,-.1){离开};
  \draw[dim,red!75!black] (0.028574763,-.65)--(.973939,-.65) node[midway,below]{最长窗邻接轮缘};
  \node[align=center] at (0,-1.95){示例 $p=0.02$，全部碰撞按三段单调支求根\\$\alpha_{\rm sup}=0.94995305$ 为擦碰相位的单侧极限};
 \end{scope}
\end{tikzpicture}

% 03 question-02
\begin{tikzpicture}
 \draw[axis] (-4.3,0)--(4.6,0) node[right]{$x$};
 \draw[axis] (0,-2.7)--(0,.9) node[above]{$y$};
 \draw[body,fill=gray!18] (-1,-1) rectangle (1,-1.35);
 \draw[blue!70!black,very thick,domain=1:4,samples=80] plot(\x,{-exp(-(\x-1)/1.4)});
 \draw[blue!70!black,very thick,domain=-4:-1,samples=80] plot(\x,{-exp((\x+1)/1.4)});
 \draw[guide] (-4,0)--(4,0);
 \draw[guide] (1,-1)--(2.4,0);\draw[guide] (1,-1)--(2.4,-1);
 \draw[->] (1.55,-1) arc(0:35.5377:.55) node[right]{$\theta_0$};
 \draw[dim] (-2,0)--(-2,-1) node[midway,left]{$l$};
 \draw[dim] (2.6,-1)--(2.6,-1.35) node[midway,right]{$h$};
 \draw[dim] (-1,-1.8)--(1,-1.8) node[midway,below]{$2r$};
 \node at (0,-2.6){边缘角取真实切线方向；曲线示意采用二维恢复模型};
\end{tikzpicture}

% 04 solution-02
\begin{tikzpicture}[x=.9cm,y=.9cm]
  \begin{scope}
    \draw[body,fill=gray!18] (-2.2,-.2) rectangle (2.2,.35);
    \draw[force,red!70!black] (0,-.2)--++(0,-1.6) node[below]{$\rho\pi r^2hg$};
    \draw[force,blue!70!black] (0,.35)--++(0,1.55) node[above]{$\rho_0g(l+h)\pi r^2$};
    \draw[force,green!50!black] (-2.2,.15)--++(-.45,1.15) node[left]{$\sigma$};
    \draw[force,green!50!black] (2.2,.15)--++(.45,1.15) node[right]{$\sigma$};
    \node at (0,-2.65){竖直分量合计 $2\pi r\sigma\sin\theta_0$};
  \end{scope}
  \begin{scope}[shift={(7.5,0)}]
    \draw[axis] (-.2,0)--(6.1,0) node[right]{$x-r$};
    \draw[axis] (0,-3.1)--(0,.8) node[above]{$y$};
    \draw[blue!70!black,very thick,domain=0:5.5,samples=100]
      plot ({\x},{-2.2*exp(-\x/1.55)});
    \draw[guide] (0,-2.2)--(5.8,-2.2) node[left]{$y(r)=-l$};
    \draw[dim] (1.55,.25)--(0,.25) node[midway,above]{$a=\sqrt{\sigma/(\rho_0g)}$};
    \node[align=center] at (3,-3.05){小坡度近似下\\$y=-l\exp[-(x-r)/a]$，$\theta_0\simeq l/a$};
  \end{scope}
\end{tikzpicture}

% 05 question-03
\begin{tikzpicture}[x=1cm,y=1cm]
  \draw[thick] (-3.8,3.1)--(1.4,3.1);
  \foreach \x in {-3.6,-3.2,...,1.2}{\draw[thin] (\x,3.1)--++(.28,.32);}
  \coordinate (O) at (-1.3,3.1); \coordinate (M) at (-3.0,-.3);
  \coordinate (D) at (-1.3,-1.0);
  \draw[very thick] (O)--(M);
  \filldraw[fill=orange!25,thick] (M) circle(.28) node[below left]{$m$};
  \draw[guide] (O)--(-1.3,-1.0);
  \draw pic[draw,->,"$\alpha_0$",angle radius=.8cm,angle eccentricity=1.35]
    {angle=M--O--D};
  \node at (-.1,.6){$\otimes\,\bm B$};
  \draw[thick] (O)--(2.55,3.1)--(2.55,1.25);
  \draw[thick] (2.05,1.25)--(3.05,1.25) (2.05,.65)--(3.05,.65);
  \draw[thick] (2.55,.65)--(2.55,-.35)--(-1.3,-.35);
  \node[right] at (3.05,.95){$C$};
  \filldraw[fill=white,thick] (-1.3,-.35) circle(.17);
  \node[align=center] at (.2,-1.2){摆通过平衡位置时短暂接通\\柔软导线使电容器充电};
\end{tikzpicture}

% 06 solution-03
\begin{tikzpicture}
 \coordinate (O) at (0,3.4);\coordinate (M) at (0,0);
 \draw[thick] (O)--(M);\fill (O) circle(2pt) node[above]{$O$};
 \draw[very thick,orange] (0,1.4)--(0,1.8) node[left]{$dx$};
 \draw[dim] (-.7,3.4)--(-.7,1.6) node[midway,left]{$x$};
 \draw[force,blue!75!black] (0,1.6)--(1.5,1.6) node[right]{$v=\omega x$};
 \node at (1.0,2.6){$\otimes\bm B$};\filldraw[fill=orange!25] (M) circle(.15) node[below]{$m$};
 \draw[force,red!75!black] (0,.7)--(-1,.7) node[left]{$dF=IB\,dx$};
 \node[draw,fill=blue!6,align=center] (a) at (5.5,3){$\mathcal E=a_B\omega$\\$a_B=Bl^2/2$};
 \node[draw,fill=orange!8,align=center] (b) at (5.5,1.4){$J(\omega_0-\omega_1)=a_Bq$\\$q=Ca_B\omega_1$};
 \node[draw,fill=green!6,align=center] (c) at (5.5,-.2){$\omega_1=\omega_0/(1+\beta)$\\$\beta=Ca_B^2/J$};
 \draw[flow] (a)--(b);\draw[flow] (b)--(c);
 \node at (2.8,-1.3){短时充电先列角冲量，再由断开后的机械能求振幅};
\end{tikzpicture}

% 07 question-04
\begin{tikzpicture}[x=1cm,y=1cm]
  \draw[body,fill=orange!8] (-2.7,-1.7) rectangle (2.7,-1.35);
  \coordinate (S) at (0,-1.35);
  \draw[force,red!70!black] (S)--++(0,2.7) node[above]{面法线};
  \draw[force,blue!70!black] (S)--++(42:3.2) node[right]{发光方向};
  \coordinate (N) at (0,1.35); \coordinate (R) at ({3.2*cos(42)},{-1.35+3.2*sin(42)});
  \draw pic[draw,->,"$\theta$",angle radius=.8cm,angle eccentricity=1.35]
    {angle=R--S--N};
  \draw[guide] (3.2,1.65) circle(.42);
  \node at (3.2,1.65){$d\Omega$};
  \node[align=center] at (0,-2.45){余弦发光面元 $dS$：\;$dE=I_0\cos\theta\,d\Omega\,dS\,dt$};
\end{tikzpicture}

% 08 solution-04
\begin{tikzpicture}[x=.8cm,y=.8cm]
  \begin{scope}
    \draw[thick] (-3.3,-1.6)--(3.3,-1.6);
    \fill (0,2.2) circle(2pt) node[above]{接收面元};
    \draw[guide] (0,2.2)--(-2.45,-1.6);
    \draw[dim] (3.0,-1.6)--(3.0,2.2) node[midway,right]{$l_0$};
    \draw[dim] (0,-2.05)--(-2.45,-2.05) node[midway,below]{$s$};
    \node[align=center] at (0,-3.0){无限平面：$\displaystyle P_1=\int_0^\infty\frac{I_0l_0^2\,2\pi s\,ds}{(s^2+l_0^2)^2}=\pi I_0$};
  \end{scope}
  \begin{scope}[shift={(8.3,.2)}]
    \draw[thick] (-3.2,2.35)--(3.2,2.35);
    \draw[body,fill=orange!10] (-2.0,2.35) ellipse[x radius=.42,y radius=.18];
    \draw[body,fill=orange!10] (2.0,2.35) ellipse[x radius=.42,y radius=.18];
    \fill (0,-1.9) circle(2pt) node[below right]{桌面点 $x$};
    \draw[guide] (0,-1.9)--(-2,2.35) (0,-1.9)--(2,2.35);
    \draw[dim] (3.0,-1.9)--(3.0,2.35) node[midway,right]{$\lambda h$};
    \draw[dim] (-2.0,2.8)--(2.0,2.8) node[midway,above]{$2h$};
    \node[align=center] at (0,-3.35){长圆柱：$P_2(d)=\pi I_0r/d$\\两灯照度在中心附近取二阶展开};
  \end{scope}
\end{tikzpicture}

% 09 question-05
\begin{tikzpicture}[x=1cm,y=1cm]
  \draw[body] (-4.8,-2.0) rectangle (-1.3,2.0);
  \draw[body] (1.3,-2.0) rectangle (4.8,2.0);
  \draw[thick,decorate,decoration={coil,aspect=.35,segment length=4pt,amplitude=3pt}]
    (-1.3,1.25)--(-.25,1.25);
  \draw[thick,decorate,decoration={coil,aspect=.35,segment length=4pt,amplitude=3pt}]
    (.25,1.25)--(1.3,1.25);
  \draw[body,fill=orange!18] (-.25,.6) rectangle (.25,1.9);
  \draw[flow] (0,.5) arc[start angle=90,end angle=-220,radius=.45] node[below]{汽轮机};
  \node[align=center] at (-3.05,0){左室：$V_0$\\初态 $p_0,T_0$\\$C_V=\dfrac32R$};
  \node[align=center] at (3.05,0){右室：$V_0$\\初态真空};
  \draw[flow] (-1.0,1.25)--(1.0,1.25) node[midway,above]{准静态输运};
  \draw[force,red!70!black] (0,-.1)--++(0,-1.15) node[below]{对外做功 $W$};
\end{tikzpicture}

% 10 solution-05
\begin{tikzpicture}[x=.86cm,y=.86cm]
  \begin{scope}
    \draw[axis] (-.2,0)--(5.5,0) node[right]{$\lambda$};
    \draw[axis] (0,-.2)--(0,4.7) node[above]{$T/T_0$};
    \draw[blue!70!black,very thick,domain=.05:1,samples=120]
      plot ({5*\x},{4*pow(\x,.4)}) node[above]{$T_L/T_0=\lambda^{2/5}$};
    \draw[guide] (5,0)--(5,4);
    \node at (2.8,-.8){刚性左室的绝热放气曲线};
  \end{scope}
  \begin{scope}[shift={(8,1.0)}]
    \node[draw,rounded corners,fill=blue!8,minimum width=3cm,minimum height=.9cm] (a) at (0,2.0){左室状态 $p_L,T_L,n_L$};
    \node[draw,rounded corners,fill=orange!10,minimum width=3cm,minimum height=.9cm] (b) at (4.5,2.0){微元气体绝热膨胀};
    \node[draw,rounded corners,fill=green!8,minimum width=3cm,minimum height=.9cm] (c) at (9,2.0){右室混合 $p_R,T_R$};
    \draw[flow] (a)--(b);
    \node[blue!70!black] at (2.25,2.72){$-dn_L$};
    \draw[flow] (b)--(c);
    \draw[force,red!70!black] (4.5,1.55)--++(0,-1.4) node[below]{$dW$};
    \node[align=center] at (4.5,-2.05){整体能量守恒确定 $p_R,T_R$\\微元第一定律给出 $W(\lambda)$ 的微分方程};
  \end{scope}
\end{tikzpicture}

% 11 question-06
\begin{tikzpicture}
 \draw[axis] (-4,0)--(4.4,0) node[right]{$z$};
 \draw[axis] (0,-1.9)--(0,2.3) node[above]{$y$};
 \draw[axis] (-1.1,-.7)--(1.4,.89) node[right]{$x$};
 \foreach \a in {-2.7,2.7}{
  \draw[body,domain=0:360,samples=90] plot({\a+.72*cos(\x)},{.46*cos(\x)+1.3*sin(\x)})--cycle;
  \draw[flow,purple,domain=15:75,samples=20] plot({\a+.72*cos(\x)},{.46*cos(\x)+1.3*sin(\x)}) node[above]{$I$};
 }
 \node[above left] at (0,0){$O$};
 \draw[dim] (0,-2.3)--(2.7,-2.3) node[midway,below]{$\lambda r$};
 \node at (0,-3.2){同向电流使轴线磁场沿 $+z$；圆心在 $(0,0,\pm\lambda r)$};
\end{tikzpicture}

% 12 solution-06
\begin{tikzpicture}[x=.85cm,y=.85cm]
 \draw[axis] (-2.6,0)--(2.8,0) node[right]{位移};
 \draw[axis] (0,-2.3)--(0,2.2) node[above]{$U-U(0)$};
 \draw[blue!75!black,very thick,domain=-2.1:2.1,samples=70] plot(\x,{.22*\x*\x});
 \draw[red!75!black,very thick,domain=-2.1:2.1,samples=70] plot(\x,{-.44*\x*\x});
 \node[blue!75!black] at (1.8,1.65){$m_0kx^2/4$};
 \node[red!75!black] at (1.6,-2.15){$-m_0kz^2/2$};
 \node[align=center] at (0,-3.4){$\lambda>1/2$：横向势阱、轴向势垒\\$\lambda=1/2$：二次项零，须算四阶};
 \begin{scope}[shift={(7.8,0)}]
  \draw[axis] (-3.1,0)--(3.6,0) node[right]{$z$};
  \foreach \a in {-2.2,2.2}{\draw[body,fill=blue!4,domain=0:360,samples=90] plot({\a+.55*cos(\x)},{.35*cos(\x)+sin(\x)})--cycle;}
  \draw[force,purple] (0,0)--(.9,0) node[above]{$\bm m$};
  \draw[force,red!75!black] (0,-.5)--(1,-.5) node[midway,below]{$\delta z$};
  \node[align=center] at (-2.2,1.9){下环 $z=-\lambda r$\\$\delta\Phi=-A\delta z$\\$\delta I=+A\delta z/(L-M)$};
  \node[align=center] at (2.2,1.9){上环 $z=+\lambda r$\\$\delta\Phi=+A\delta z$\\$\delta I=-A\delta z/(L-M)$};
  \draw[force,green!50!black] (.8,-1.2)--(-.4,-1.2) node[midway,below]{$F_{\rm ind}=-2A^2\delta z/(L-M)$};
 \end{scope}
\end{tikzpicture}

% 13 question-07
\begin{tikzpicture}[x=1cm,y=1cm]
  \draw[axis] (-4.5,0)--(4.8,0) node[right]{$x$};
  \draw[axis] (0,-3.2)--(0,5.1) node[above]{$y$};
  \draw[body,fill=blue!8] (-3,0) arc[start angle=180,end angle=0,radius=3]--cycle;
  \foreach \x in {-2.4,-1.6,-.8,0,.8,1.6,2.4}{
    \draw[flow,orange!80!black] (\x,-2.8)--(\x,-.05);}
  \draw[dim] (0,.28)--(3,.28) node[midway,above]{$r$};
  \node at (0,1.1){$n=1.56$};
  \node[align=center] at (0,-3.6){平行光由平面入射半圆柱透镜\\只在弧面发生折射};
\end{tikzpicture}

% 14 solution-07
\begin{tikzpicture}[x=.78cm,y=.78cm]
  \draw[axis] (-7.0,0)--(5.5,0) node[right]{$x$};
  \draw[axis] (0,-1.2)--(0,10.25) node[above]{$y$};
  \draw[body,fill=blue!7] (-3.4,0) arc[start angle=180,end angle=0,radius=3.4]--cycle;
  \coordinate (V) at (0,3.4); \coordinate (F) at (0,9.4714);
  \fill (F) circle(2pt) node[right]{$F(0,h_0)$};
  \draw[flow,orange!80!black] (.25,-1)--(.25,3.3908)--(F);\draw[guide] (0,3.4)--(F);
  \coordinate (C) at ({3.4/1.56},{sqrt(3.4^2-(3.4/1.56)^2)});
  \draw[flow,red!75!black] ({3.4/1.56},-1)--(C)--(-6.036,9.4714);
  \draw[guide] (0,0)--(C);
  \draw[guide] (-6.45,9.4714)--(4.05,9.4714) node[right]{焦点平面};
  \draw[guide] (-4.0,6.8)--(4.05,6.8) node[right]{$y=10.0\,\mathrm{cm}$};
  \draw[dim] (4.5,0)--(4.5,9.4714) node[midway,right]{$h_0=r+\dfrac{r}{n-1}$};
  \draw[dim] (0,-.65)--({3.4/1.56},-.65) node[midway,below]{$\xi_{\max}=r/n$};
  \node[align=center] at (0,-2.35){临界光线在弧面处沿切线出射；\\它与近轴焦点共同决定球差宽度和照亮边界};
\end{tikzpicture}

% 15 question-08
\begin{tikzpicture}
 \draw[thick] (-4,0)--(4,0);
 \draw[body,fill=blue!6] (-3.2,3.2) arc(180:360:3.2)--cycle;
 \coordinate (P) at ({-3.2*sqrt(5)/3},{3.2/3});
 \fill[orange!80!black] (P) circle(2pt);
 \draw[guide] (0,3.2)--(P);\draw[guide] (-3.6,{3.2/3})--(P);
 \draw[dim] (-3.6,0)--(-3.6,{3.2/3}) node[midway,left]{$R/3$};
 \draw[dim] (0,3.2)--(3.2,3.2) node[midway,above]{$R$};
 \draw[force,blue] (P)--++({-1.2*2/3},{1.2*sqrt(5)/3}) node[left]{子午切向};
 \node[right,red!75!black] at (P){$\odot v_{\varphi0}$};
 \node[align=center] at (0,-.8){初始点满足 $\rho_0=\sqrt5R/3$\\水平切向速度垂直于子午截面};
\end{tikzpicture}

% 16 solution-08
\begin{tikzpicture}[x=.85cm,y=.85cm]
 \draw[body,fill=blue!6] (-3,3) arc(180:360:3)--cycle;
 \coordinate (P) at ({-sqrt(5)},1);\fill[orange!80!black] (P) circle(2pt);
 \draw[force,blue] (P)--++({-.95*2/3},{.95*sqrt(5)/3}) node[above left]{向上};
 \draw[force,green!50!black] (P)--++({.95*2/3},{-.95*sqrt(5)/3}) node[below right]{向下};
 \draw[flow,blue] (-2.8,1.96) arc(200:184:3);
 \node[align=center] at (0,-1.1){两种 $L_z=0$ 的子午方向\\到沿阈值 $v^2=4gR/3$，实际越出取严格大于};
 \begin{scope}[shift={(7.5,0)},x=4cm,y=1.5cm]
  \draw[axis] (.08,0)--(1.12,0) node[right]{$h/R$};
  \draw[axis] (.1,-.1)--(.1,3.7) node[above]{$U_{\rm eff}/mgR$};
  \draw[red!75!black,very thick,domain=.14:1,samples=151] plot(\x,{\x+(5/6)/(2*\x-\x*\x)});
  \draw[guide] (1,0)--(1,1.833333);\draw[guide] (.333333,0)--(.333333,1.833333);
  \node[below] at (.333333,0){$1/3$};\node[below] at (1,0){$1$};
  \draw[guide] (.15,1.833333)--(1,1.833333) node[right]{$E/mgR=11/6$};
  \node[align=center] at (.6,-.7){水平发射阈值的真实曲线\\$u+(5/6)/(2u-u^2)$};
 \end{scope}
\end{tikzpicture}

\end{document}
